\section{Symbolic parity and the congruence tower}
\label{sec:parity}

The first local correction has an exact symbolic description.  Reducing the
generators modulo two gives
\[
 J:=\overline A=\begin{pmatrix}1&1\\1&1\end{pmatrix},
 \qquad
 E:=\overline B=\begin{pmatrix}1&0\\0&0\end{pmatrix},
 \tag{4.1}\label{4.1}
\]
with
\[
 J^2=0,\qquad E^2=E,\qquad (EJ)^2=EJ,
 \qquad \tr(EJ)=1.
 \tag{4.2}\label{4.2}
\]

\begin{theorem}[Cyclic golden-mean gate]
\label{thm:parity}
For every based cyclic word \(w\),
\[
 2\mid D_w
 \quad\Longleftrightarrow\quad
 w\text{ has no cyclic occurrence of \texttt{aa}}.
 \tag{4.3}\label{4.3}
\]
The number of such based words of length \(n\) is the Lucas number
\[
 L_n=\tr\begin{pmatrix}0&1\\1&1\end{pmatrix}^{n}
 =\golden^n+(-\golden^{-1})^n.
 \tag{4.4}\label{4.4}
\]
\end{theorem}

\begin{proof}
Because \(n\ge1\), equation \eqref{3.4} gives
\(D_w\equiv1+\tr M_w\pmod2\).  If the cyclic word contains
\texttt{aa}, rotate the trace until \(J^2\) is adjacent; its trace is then
zero.  If the word has no cyclic \texttt{aa}, it is either all \texttt{b},
whose reduction is \(E\), or every \(J\) is separated by a nonempty
\(E\)-run.  Idempotence compresses the cyclic trace to
\(\tr((EJ)^k)=1\).  This proves \eqref{4.3}.  Equation \eqref{4.4} is the
closed-walk count for the two-state no-\texttt{aa} graph.
\end{proof}

The higher dyadic layers remain chronological.  For any positive integer
\(d\),
\[
 \oddpart(d)
 =d\left(1-\sum_{k\ge1}2^{-k}\mathbf 1_{2^k\mid d}\right).
 \tag{4.5}\label{4.5}
\]
Indeed, if \(\vTwo(d)=m\), the finite sum in parentheses equals \(2^{-m}\).
Define
\[
 C_{n,k}=\sum_{\substack{|w|=n\\2^k\mid D_w}}D_w.
 \tag{4.6}\label{4.6}
\]
Then the complete fixed count is the coefficientwise finite identity
\[
 \Nsol_n=\Ninf_n-\sum_{k\ge1}2^{-k}C_{n,k}.
 \tag{4.7}\label{4.7}
\]

\begin{proposition}[Finite chronological recurrence at fixed depth]
\label{prop:tower}
For every fixed \(k\), the sequence \(C_{n,k}\) is generated by a finite
linear recurrence on chronological residue states modulo \(2^k\).
\end{proposition}

\begin{proof}
For a reachable residue matrix \(g\in M_2(\Z/2^k\Z)\), retain the word
count \(c_n(g)\), the scaled count \(q_n(g)=8^n c_n(g)\), and the integer
matrix sum
\[
 S_n(g)=\sum_{\substack{|w|=n\\M_w\equiv g\ (2^k)}}M_w.
\]
Appending the symbol with matrix \(L\in\{A,B\}\) on the right of the word
multiplies the return on the left.  Thus
\[
 c_{n+1}(h)=\sum_{Lg=h}c_n(g),\qquad
 q_{n+1}(h)=8\sum_{Lg=h}q_n(g),\qquad
 S_{n+1}(h)=\sum_{Lg=h}L S_n(g).
 \tag{4.8}\label{4.8}
\]
The residue \(8^n\bmod2^k\) may be included as one more finite state.  The
selector \(2^k\mid D_w\) then depends only on that state and \(g\), and
\[
 C_{n,k}=\sum_{\substack{g:\;1+8^n-\tr g\equiv0\ (2^k)}}
 \bigl(c_n(g)+q_n(g)-\tr S_n(g)\bigr).
 \tag{4.9}\label{4.9}
\]
This is a finite linear system for each fixed \(k\).
\end{proof}

Proposition~\ref{prop:tower} does not construct one finite transfer matrix
for the exact odd-part weight.  Equation \eqref{4.7} uses unbounded congruence
depth, is exact one coefficient at a time, and supplies no license to
interchange an infinite state sum with analytic continuation.

\begin{table}[t]
\centering
\small
\caption{Exact counts.  The active-word column is the Lucas number from
Theorem~\ref{thm:parity}.}
\label{tab:counts}
\begin{tabular}{r@{\quad}r@{\quad}r@{\quad}r@{\quad}r}
\toprule
\(n\)&\(\Ninf_n\)&\(\Nsol_n\)&\(\Delta_n\)&\(L_n\)\\
\midrule
1&5&4&1&1\\
2&157&56&101&3\\
3&3194&2446&748&4\\
4&57121&34504&22617&7\\
5&969035&807804&161231&11\\
6&16021054&11965154&4055900&18\\
7&261230555&232113186&29117369&29\\
8&4226257633&3511221904&715035729&47\\
9&68064014258&63051464302&5012549956&76\\
10&1093258037257&972315505846&120942531411&123\\
\bottomrule
\end{tabular}
\end{table}
